\documentclass[11pt,a4paper]{article}

% ---------- Packages ----------
\usepackage[margin=2.5cm]{geometry}
\usepackage{amsmath, amssymb, amsthm, mathtools}
\usepackage{enumitem}
\usepackage[hidelinks]{hyperref}
\usepackage{fancyhdr}

% ---------- Sheet details (edit these) ----------
\newcommand{\modulename}{Introduction to Pure Mathematics}
\newcommand{\sheetnumber}{1}
\newcommand{\myname}{Daniel Greenfield}

% ---------- Header ----------
\setlength{\headheight}{14pt}
\pagestyle{fancy}
\fancyhf{}
\lhead{\modulename}
\rhead{Sheet \sheetnumber}
\cfoot{\thepage}

% ---------- Number sets ----------
\newcommand{\N}{\mathbb{N}}
\newcommand{\Z}{\mathbb{Z}}
\newcommand{\Q}{\mathbb{Q}}
\newcommand{\R}{\mathbb{R}}
\newcommand{\C}{\mathbb{C}}




\begin{document}

\begin{center}
  {\Large \textbf{\modulename}}\\[0.3em]
  {\large Problem Sheet \sheetnumber}\\[0.3em]
  \myname
\end{center}
\section*{Question 1.1.}

    \begin{enumerate}[label=\alph*.]
        \item  Let $u = \begin{pmatrix} 1 \\ 0 \end{pmatrix}$ and $v = \begin{pmatrix} 1 \\ y\end{pmatrix}$ with $y  \in \R$ then 

            $ \| u \| = 1$

            $\| v \| = \sqrt{1 + y^2}$

            $\| u + v \| = \sqrt{4 + y^2}$

            Finding $ y \in \R $ such that $\| u + v\| = \| u \| + \| v \| $.

            \[
                \sqrt{4 + y^2} = 1 + \sqrt{1 + y^2}
            \]
            
            Squaring both sides gives us
            
            \[
                4 + y^2 = 1 + 1 + y^2 + 2 \cdot \sqrt{1 + y^2}
            \]

            Which can be simplified to:

            \[
                1 = \sqrt{1 + y^2}
            \]

            Squaring both sides again leaves us with
            
            \[
                1 = 1 + y^2
            \]

            Therefore $ y = 0 $.
        \item let $v_1, v_2, \cdots , v_k$ be k arbitrary vectors in $\R^n$. For $k \in \Z_{+}$ let $P(k)$  be the statement:
            \[
                \| v_1 + v_2 + \cdots + v_k \| \leq \| v_1 \| + \| v_2 \| + \cdots \| v_k \|
            \]
        We will use induction to prove that $\forall k \in \Z_+ P(k)$ is true.
        For P(1) \[
            \| v_1 \| \leq \| v_1 \| 
        \] 
        is true as the left hand side and the right hand side are equal. \\
        For P(2) \[
            \| v_1 + v_2 \| \leq \| v_1 \| + \| v_2 \|      
        \] is true due to the triangle inequality ($ \| u + v \| \leq \| u \| + \| v \| $).
        Now assume that $P(k)$ is true for some integers $k > 2$.
        \[
            \| v_1 + v_2 + \cdots + v_k \| \leq \| v_1 \| + \| v_2 \| + \cdots \| v_k \|
        \]
        Adding $\| v_{k+1} \|$ where $v_{k+1} \in \R^n$ is some arbitrary vector in $\R^n$ to both sides will not change the truth value of the inequality giving us
        \[
            \| v_1 + v_2 + \cdots + v_k \| + \| v_{k+1} \|  \leq \| v_1 \| + \| v_2 \| + \cdots \| v_k \| + \| v_{k+1} \| 
        \]
        Using the triangle inequality again we can get a second inequality involving the left hand side:
        \[
            \| v_1 + v_2 + \cdots + v_k + v_{k+1}\|\leq \| v_1 + v_2 + \cdots + v_k \| + \| v_{k+1} \|  \leq \| v_1 \| + \| v_2 \| + \cdots + \| v_k \| + \| v_{k+1} \| 
        \]
        Due to the transitive properties of inequalities we get that:
        \[
            \| v_1 + v_2 + \cdots + v_k + v_{k+1}\| \leq \| v_1 \| + \| v_2 \| + \cdots + \| v_k \| + \| v_{k+1} \| 
        \]
        Which is the statement $P(k+1)$, so we have shown that if $P(k)$ is true then $P(k+1)$ is true.

        So we have shown that for $k > 2$, $P(k) \Rightarrow P(k+1)$ and that $P(1)$ and $P(2)$ are true. Then by the method of induction we have shown that $P(k)$ is true for all $k \in \Z_+$.
       
        An example for which $\|v_1 + v_2 + \cdots + v_k\| = \| v_1 \| + \| v_2 \| + \cdot + \| v_k \|$  is \\ $v_1 = \begin{pmatrix} 1 \\ 1 \\ -4 \end{pmatrix}, \ \ v_2 = \begin{pmatrix} 2 \\ 2 \\ -8 \end{pmatrix} , \ \ v_3 = \begin{pmatrix}3 \\ 3 \\ -12\end{pmatrix}$. 
        \\ Where $\| v_1\| = 3\sqrt{2}$ , $\| v_2\| = 6\sqrt{2}$, $\| v_3 \| = 9\sqrt{2}$, $v_1 + v_2 + v_3 = \begin{pmatrix}6 \\ 6 \\ -24\end{pmatrix}, \\ 
        $\|v_1 + v_2 + v_3\| = 18\sqrt{2}$ and $\| v_1 \| + \| v_2 \| + \| v_3 \| = 18\sqrt{2} = \| v_1 + v_2 + v_3 \|$
        \\ However this equality holds for any set of vectors if all the vectors are parallel in the same direction ($\forall i,j \in [1,k], \exists \lambda \in \R_+ $ such that $v_i = \lambda v_j).



    \end{enumerate}



\end{document}
